% ═══════════════════════════════════════════════════════════════════════════════
\section{Discussion: Synthesis, Theory Revisited, and the Road Not Taken}
\label{sec:discussion}
% ═══════════════════════════════════════════════════════════════════════════════

\subsection{Revisiting the Central Research Question}

This dissertation asked: \textit{to what extent can Uzbekistan achieve structural transformation through strategic integration into Northeast Asian production networks?} The empirical evidence assembled across Sections~4--12 allows a nuanced, graduated answer rather than a binary one.

\begin{table}[H]
\centering
\caption{Synthesis of Empirical Findings and Implications}
\label{tab:synthesis_matrix}
\resizebox{\textwidth}{!}{
\begin{tabular}{p{0.25\linewidth}p{0.45\linewidth}p{0.3\linewidth}}
\toprule
\textbf{Evidence Base / Method} & \textbf{Key Empirical Finding} & \textbf{Implication for GVC Governance} \\
\midrule
\textbf{Chow Structural Break Test} & 2017 confirmed as a structural break for manufacturing ($F=6.937, p=0.013$) & Reforms successfully accelerated industrial output but without necessarily altering GVC positioning. \\
\addlinespace
\textbf{RCA Index Analysis} & Comparative advantages remain concentrated in resources/low-complexity goods & Failure to achieve functional upgrading; lock-in to primary/extractive nodes. \\
\addlinespace
\textbf{OLS \& Panel Regression} & Intermediate goods imports significant ($\beta=2.34$); FDI volume insignificant ($\beta=-0.014$) & Upgrading depends on network embedding (components) rather than sheer capital volume. \\
\addlinespace
\textbf{Trade Network Metrics} & High betweenness (0.26) and in-degree; low forward integration & Conduit for imported components but lacks reciprocal high-value exports. \\
\addlinespace
\textbf{Comparative Benchmarking} & FDI reversed post-reform (4.1\% $\to$ 2.4\%), unlike Vietnam & Fails to attract the export-platform FDI needed for relational governance. \\
\addlinespace
\textbf{Constraint Diagnostics} & Logistics cost and transit time are the absolute binding constraints & Structural barrier preventing transition from 'captive' assembly to 'relational' integration. \\
\bottomrule
\end{tabular}}
\end{table}

This synthesis (Table~\ref{tab:synthesis_matrix}) draws together the methodological components to demonstrate the central argument: Uzbekistan has deepened its integration into Northeast Asian networks, but remains trapped in a \textit{captive} governance structure where it acts as a passive absorber of intermediate goods rather than an active relational exporter.
The affirmative case is real and documented. The 2017 Mirziyoyev reforms constituted a genuine structural break in Uzbekistan's manufacturing trajectory (Chow $F=6.937$, $p=0.013$), accelerating annual manufacturing VA growth from 0.4 to 1.2 percentage points. Northeast Asian trade integration has deepened rapidly and at scale: bilateral trade with China increased 280\% between 2017 and 2023, and the intermediate goods share of imports from all three Northeast Asian partners (China: 85\%, South Korea: 90\%, Japan: 98\%) is fully consistent with GVC backward integration potential. The ECI has improved measurably ($-0.65 \to -0.2$), and Uzbekistan's reform trajectory in manufacturing VA at years 0--5 post-reform broadly matches Vietnam's equivalent trajectory.

The critical qualifications are equally documented. Uzbekistan's FDI trajectory has \textit{reversed} post-reform, declining from 4.1\% of GDP in 2017 to 2.4\% in 2023, while Vietnam and Thailand both saw FDI rise continuously through comparable reform windows. The dominant FDI source (China, 28\%) is primarily infrastructure-linked rather than export-manufacturing-linked. Japan and South Korea --- the partners whose capital most closely resembles the export-platform FDI that drove Vietnamese upgrading --- remain marginal FDI contributors despite their substantial trade relationships. Export sophistication, measured by ECI and manufacturing RCA, has improved only modestly and remains below negative-zero on the global distribution. The GVC positioning analysis (Figure~\ref{fig:gvc_positioning}) reveals a telling trajectory: Uzbekistan has moved rightward (high backward integration) but not upward (low forward integration), the quantitative signature of \textit{captive} rather than relational GVC governance.

The calibrated answer is therefore: \textit{Uzbekistan can achieve significant structural transformation through Northeast Asian GVC integration, but only if the current trajectory --- which favours infrastructure-linked Chinese capital over manufacturing-linked Japanese and Korean capital --- is deliberately redirected, and the binding constraints of logistics inefficiency and governance quality are systematically reduced. Without these interventions, the most likely scenario is import-led consumption growth and resource-processing export growth, rather than the export-sophistication path that would generate sustained income convergence.}

\subsection{What the Regression Results Actually Say: The Mechanism Is Intermediate Goods, Not FDI Volume}

The econometric evidence presented in Section~\ref{sec:regression} carries a policy implication that is sharper than the aggregate discussion of investment promotion typically acknowledges. The OLS results reveal that the coefficient on lagged FDI as a share of GDP is negative but statistically insignificant ($\beta = -0.014$, $p > 0.10$), while the coefficient on the intermediate goods import share is positive and significant ($\beta = 2.34$, $p < 0.05$). This juxtaposition is not a marginal nuance --- it is the central empirical finding, and it directly challenges the generic ``attract more FDI'' logic that dominates Uzbekistan's investment promotion strategy.

The finding implies that the operative mechanism linking Northeast Asian economic engagement to manufacturing upgrading in Uzbekistan is not the \textit{volume} of capital inflows, but the \textit{composition} of trade relationships --- specifically, the depth of intermediate goods integration. Countries and firms that embed themselves in production networks as processors of imported components (backward integration) generate the production experience, quality standards, and supplier relationships that eventually translate into forward integration and export sophistication. FDI, in this framework, matters primarily to the extent that it \textit{generates} intermediate goods trade relationships rather than as an independent stimulus. Infrastructure-linked FDI that does not activate intermediate goods trade --- the dominant Chinese BRI model in Uzbekistan --- will therefore produce near-zero manufacturing upgrading effects regardless of its headline value.

The policy implication is precise: first-order instruments are those that deepen intermediate goods trade relationships --- supplier development programmes, backward linkage requirements written into SEZ operating agreements, joint venture structures that mandate domestic component sourcing thresholds. Second-order instruments are those targeting FDI quantity (tax holidays, land subsidies, administrative streamlining). Both matter, but the regression evidence indicates that the current policy architecture inverts the priority ordering.

\subsection{The Flying Geese Mechanism: Operating But Distorted}

The flying geese mechanism, as formalised by \citet{akamatsu1962} and updated for modern GVC analysis, predicts that as lead economies progressively vacate labour-intensive production tasks, follower economies with abundant cheap labour and improving institutions should attract the relocating capital. The evidence confirms that this mechanism is operating in Uzbekistan's direction: Chinese factory relocation from coastal provinces, South Korean \textit{China+1} strategies, and Japanese supply chain diversification post-COVID are all directing capital toward Central Asia. However, as noted in the development literature, global economic evolution patterns often bypass landlocked interiors unless proactive state intervention alters the cost-benefit calculus.

The distortion is not merely a stall but a fundamental divergence in the \textit{type} of capital flowing relative to model predictions. The flying geese framework anticipates export-manufacturing-linked FDI (the Japanese/Korean model) seeking wage arbitrage. Instead, Uzbekistan is disproportionately absorbing infrastructure-linked Chinese FDI. This capital builds roads and power plants but does not directly integrate domestic firms into manufacturing export networks, explaining the high backward integration (capital imports) without commensurate forward integration (manufactured exports).

However, this divergence may represent a necessary sequencing adaptation rather than a terminal failure. As established in the constraint hierarchy (Section~\ref{sec:constraints}), landlockedness and energy reliability are the absolute binding constraints on export competitiveness. Chinese infrastructure FDI may therefore serve as a prerequisite enabling condition---building the physical connectivity and power grid stability required before Japanese and Korean electronics or automotive manufacturers will commit the export-platform FDI necessary to complete the flying geese cycle.

The flying geese literature's most relevant recent contribution is the concept of ``GVC participation without upgrading'' --- the documented tendency of developing countries in the post-2000 era to integrate into GVCs (increasing backward integration) without the domestic capability accumulation (technology adoption, product diversification, supplier development) that earlier theoretical models predicted would follow \citep{milberg2013}. Their empirical analysis shows that GVC participation is positively associated with export growth but only weakly associated with productivity growth or ECI improvement --- precisely the pattern observed in Uzbekistan. This suggests that the stall at flying geese stage one is not peculiar to Uzbekistan but is a systemic feature of the contemporary GVC landscape that requires active policy rather than passive openness.

\subsection{The Vietnam Comparison: What Uzbekistan Has and Has Not Replicated}

The dissertation's most detailed comparator, Vietnam, achieved its GVC integration breakthrough not through geography (Vietnam's coastal access is an advantage, not a constant), but through the \textit{combination} of three factors that operated simultaneously: (i) a credible trade policy commitment locked in through WTO accession and the US-Vietnam BTA; (ii) a specific set of anchor investors (Samsung at Thai Nguyen, LG at Hai Phong, Intel at HCMC) whose scale of investment created an entire ecosystem of Tier 2--3 suppliers requiring domestic sourcing; and (iii) a deliberate human capital development strategy calibrated to the specific skills demands of electronics manufacturing (STEM education expansion, vocational training in collaboration with Samsung and LG).

Of these three factors, Uzbekistan currently has partial versions of only the first and none of the second. Trade policy credibility is being built through the WTO accession process and through the 2022 GSP+ application to the EU, but remains incomplete relative to Vietnam's locked-in commitments. Anchor manufacturing investors are absent: no Japanese or Korean firm has committed investment at the Samsung-in-Vietnam scale that would create a domestic supplier ecosystem. Human capital development is improving but is not yet calibrated to the specific demands of Northeast Asian manufacturing partners.

The implication is not that Uzbekistan cannot follow the Vietnamese path, but that the sequencing matters. Vietnam's FDI surge was not a response to abstract policy openness --- it was triggered by Samsung's investment decision in 2008, which was itself triggered by Samsung's specific business logic (diversifying Samsung SDI battery supply chains away from Chinese single-sourcing). Uzbekistan needs an equivalent anchor investment decision from a Northeast Asian industrial giant to trigger the ecosystem effects that generalised investment promotion cannot achieve.

\subsection{The Thailand Comparison: Sectoral Targeting over Spatial Concentration}

Thailand's Eastern Seaboard model adds a dimension that the Vietnam comparison obscures: the importance of \textit{sectoral targeting} in SEZ strategy. Thailand did not simply designate land and offer tax holidays; it identified a specific target sector (automotive), committed infrastructure investment calibrated to that sector (deep-sea port for vehicle export, specialised steel strip for body panels, industrial water supply for paint processes), and offered sufficiently differentiated incentives to attract Tier 1 anchor investors (Toyota, Honda) whose presence then pulled Tier 2 and Tier 3 suppliers from Japan to Thailand.

Uzbekistan's current SEZ landscape lacks this sectoral focus: zones are designated across eight sectors (electronics, textiles, chemicals, metals, logistics, ICT, food processing, construction materials) without sufficient infrastructure differentiation across zones to make any one zone genuinely world-class for a specific sector. The Navoiy aviation-logistics model represents a partial exception: it is meaningfully differentiated and has attracted an anchor investor (Korean Air). The analytical prescription from the Thailand comparison is to replicate the Navoiy model's sectoral focus in at least two additional zones, targeting the automotive and electronics sectors where Northeast Asian partners have the deepest interest and Uzbekistan has the clearest labour cost advantage.

\subsection{WTO Accession as Institutional Lock-In: The Vietnam Parallel}

The comparative and structural analyses converge on a critical institutional vulnerability: the reversibility of decree-based reform. As identified in Section~\ref{sec:constraints}, the political economy of Uzbekistan's liberalisation relies heavily on executive initiative rather than deeply embedded institutional independence. This creates a credibility gap for Northeast Asian investors weighing long-term, capital-intensive manufacturing commitments. Vietnam resolved an identical credibility gap not through domestic institutional maturation alone, but through external lock-in. Its 2007 WTO accession provided an irreversible commitment mechanism that assured foreign investors that tariff reductions, intellectual property protections, and dispute resolution frameworks would survive domestic political transitions. For Uzbekistan, WTO membership serves the exact same function. It is not merely a trade facilitation exercise, but the definitive signal that the 2017 liberalisation trajectory cannot be unwound by future administrations. This mechanism directly underpins Policy Implication 5: accelerating WTO accession is the most potent available instrument to signal the permanent transition from a closed, state-directed economy to a rule-bound node in global production networks.

\subsection{The Premature Deindustrialisation Counter-Thesis and Why It Does Not Foreclose the Uzbek Case}

\citet{rodrik2016} argues that the manufacturing-led convergence path that worked for Korea, Taiwan, and Vietnam is now structurally narrowing for late developers. Labour-saving automation is reducing the employment content of manufacturing, and the shift of GVC value toward pre- and post-production segments means that assembly-stage entry yields diminishing domestic value-added. \citet{diao2019} extend this critique to show that structural change in many developing countries is moving labour from agriculture to services rather than manufacturing, failing to replicate the productivity convergence the East Asian trajectory achieved.

This critique is taken seriously here but does not invalidate the Uzbekistan case for the following reasons. First, Uzbekistan is not competing primarily for the automatable, low-skill, high-volume assembly tasks that the premature deindustrialisation literature identifies as most vulnerable. Its demographic profile, STEM education base, and EAEU market access position it to compete for medium-skill, relationship-intensive manufacturing (automotive components, pharmaceuticals, precision instruments) where automation pressures are lower and relationship governance --- the relational GVC tier --- remains labour-intensive at the engineer and technician level. Second, the premature deindustrialisation thesis is most acute for large-population, low-income economies facing mass employment absorption demands (South Asian and sub-Saharan African LDCs). For a middle-income, urbanised economy of Uzbekistan's profile, the critical constraint is institutional and logistical quality, not the structural foreclosure of the manufacturing pathway. Third, the regional market dimension is absent from Rodrik's aggregate global argument: Uzbekistan's proximity to a 180 million-person EAEU market, where trade preferences substitute partially for the comparative advantage that coastal location provides in global maritime trade, creates a viable mid-tier manufacturing niche that the premature deindustrialisation thesis does not capture.

In short, the flying geese mechanism is distorted but not foreclosed for Uzbekistan. The critical question is whether the current window --- created by China+1 supply chain diversification and post-COVID resilience reconfiguration --- is seized with anchor investor targeting and institutional credibility, or allowed to close under the weight of logistics bottlenecks and governance gaps.

\subsection{The Limits of This Analysis}

Academic rigour requires explicit acknowledgement of what this dissertation cannot establish. The OLS regressions establish associations, not causal effects. The small sample ($N=14$) means that coefficients are imprecisely estimated and F-statistics have limited power against plausible alternatives. The comparative benchmarking normalises for reform-stage but not for initial conditions (population, coastline, colonial history, regional power structure). Furthermore, relying on aggregate ECI data masks intra-sectoral upgrading dynamics that might be occurring at the firm level.

More fundamentally, this analysis describes the conditions for GVC integration at a point in time (2017--2023) that is itself a moment of global supply chain restructuring of unusual intensity: COVID-19 disruptions, US-China trade tensions, and global supply chain reconfigurations are reshaping global production networks in ways that may create conditions for Central Asian integration that did not exist a decade ago. The opportunity window identified in this dissertation may be more time-sensitive than the long-run structural analysis suggests.

Future research priorities include: panel regression analysis using a comparator group of Central Asian and landlocked reforming economies; firm-level survey data on technology transfer within Uzbekistan's SEZs to directly test the captive/relational governance distinction \citep{gereffi2005}; and forward-looking input-output analysis using OECD TiVA data to quantify the domestic value-added content of Uzbekistan's manufacturing output.
